% Generated draft: requires mathematical and source review.
\documentclass[11pt]{article}
\usepackage[margin=25mm]{geometry}
\usepackage{fontspec}
\setmainfont{Latin Modern Roman}
\usepackage{amsmath,amssymb,mathtools,unicode-math}
\setmathfont{Latin Modern Math}
\usepackage{enumitem,needspace,xurl,hyperref}
\setlength{\parindent}{0pt}
\setlength{\parskip}{5pt}
\setlength{\emergencystretch}{3em}
\newcommand{\N}{\mathbb{N}}
\newcommand{\Z}{\mathbb{Z}}
\newcommand{\Q}{\mathbb{Q}}
\newcommand{\R}{\mathbb{R}}
\newcommand{\C}{\mathbb{C}}
\newcommand{\F}{\mathbb{F}}
\newcommand{\A}{\mathbb{A}}
\newcommand{\PP}{\mathbb{P}}
\newcommand{\E}{\mathbb{E}}
\newcommand{\eps}{\varepsilon}
\newcommand{\dd}{\,\mathrm{d}}
\begin{document}
\section*{UnsolvedMath ID: 701414}
\textbf{Source problem code:} catalogue entry
\subsection{Definitions and mathematical statement}
Let \(p\) be a prime. A finite \(p\)-group is a finite group of order \(p^n\) for some integer \(n\geq 0\). Say that \(G\) is two-generated if \(G=\langle x,y\rangle\) for some \(x,y\in G\). Define its lower central series by \(\gamma_1(G)=G\) and \(\gamma_{i+1}(G)=[\gamma_i(G),G]\), where \([H,G]\) is generated by the commutators \(h^{-1}g^{-1}hg\) with \(h\in H\) and \(g\in G\). For nontrivial nilpotent \(G\), its nilpotency class \(c(G)\) is the integer \(c\geq 1\) such that \(\gamma_c(G)\neq\{1\}\) and \(\gamma_{c+1}(G)=\{1\}\).

Under the standard presentation-theoretic convention, the deficiency of a finitely presented group is
\[
\operatorname{def}(G)=\sup\bigl\{|X|-|R|:G\cong\langle X\mid R\rangle,\ X\text{ and }R\text{ finite}\bigr\}.
\]
The source refers to Problem 8.12 for its meaning of deficiency, but that referenced text is not supplied. Consequently, use of this definition here is conditional on agreement with that reference. Two-generation is a property of \(G\); this formulation does not additionally require that a presentation attaining deficiency zero use two generators.

Interpreting the question as concerning each fixed prime \(p>3\), J. Wiegold asks whether
\[
\forall\,p\text{ prime with }p>3\quad
\forall\,C\in\Z_{\geq 1}\quad
\exists\,G:\quad
\begin{cases}
G\text{ is a finite }p\text{-group},\\
G=\langle x,y\rangle\text{ for some }x,y\in G,\\
\operatorname{def}(G)=0,\\
c(G)\geq C.
\end{cases}
\]
The supplied wording does not explicitly settle whether the prime must be fixed while the class grows. If primes may vary with the class bound, the quantifiers instead become \(\forall C\in\Z_{\geq 1}\ \exists p>3\text{ prime}\ \exists G\), with the same conditions on \(G\). These readings are retained as an unresolved qualification. If ``two-generator'' means minimum generating number exactly two, imposing \(C\geq 2\) excludes cyclic groups and leaves the unbounded-class question unchanged.
\subsection{Short English statement}
For each fixed prime \(p>3\), can finite groups generated by two elements, of order a power of \(p\), and of deficiency zero have arbitrarily large nilpotency class? The referenced definition of deficiency and the intended quantification over primes require confirmation.
\subsection{Sources}
\begin{itemize}
\item Local UnsolvedMath registry: supplied complete source record, permanent ID \(701414\), exact source code \texttt{catalogue entry}, titled \emph{Kourovka Notebook Problem 17.113}.
\item Supplied attribution: J. Wiegold, proposed in 2010; original problem number 17.113 in \emph{The Kourovka Notebook (Unsolved Problems in Group Theory)}, 17th issue (2010).
\item The supplied background also identifies the 21st issue edition (2026), \texttt{arXiv:1401.0300v48}. The collection description names E. I. Khukhro and V. D. Mazurov as editors. These bibliographic details are supplied attribution and have not been independently checked.
\item The statement explicitly refers to Problem 8.12 for deficiency; its contents are absent from the supplied record.
\item The registry labels the problem \emph{open} and records an update timestamp of \texttt{2026-10-16T14:00:00Z}. This is supplied registry metadata, not an independently verified assessment of the literature.
\end{itemize}
\end{document}
